%% parts/p07/lambda-calculi-and-type-theories.tex --- Phase-4 P7 chapter 7.4
%% New-content/depth deliverables: D-P9-003, D-P9-004, N-a-015,
%% N-c-029, N-c-030, and apparatus A-073.  Pointer-only boundaries:
%% N-x-001, N-x-002, N-x-011, N-x-012.
%% This chapter is new material: it has no source-compression denominator.
%% Rules remain in registry/calculi; the body contains statements and lookup
%% apparatus only.

\chapter{Lambda Calculi and Type Theories}
\label{sec:p07-lambda-calculi-and-type-theories}
\chaptermark{Lambda Calculi and Type Theories}

This chapter separates four surfaces that are too often conflated: untyped
conversion, simply typed terms, explicitly polymorphic System~F, and an
explicitly bounded intensional dependent-type entry.  Each metatheorem is scoped to its stated
syntax.  In particular, confluence is not normalization, Church-style
checking is not Curry-style typability, and ``Martin--L\"of type theory'' is
not treated as the name of one immutable calculus.  Standard sources are
\cite{Church1936,ChurchRosser1936,Barendregt1984,
GirardLafontTaylor1989,MartinLof1984}.

\section{Untyped terms, substitution, and conversion}

\begin{hbdefinition}{def:p07-untyped-lambda-terms}{Untyped lambda terms}
Fix a countably infinite set of term variables.  The \emph{untyped lambda
terms} form the least set that contains every term variable and is closed
under application and abstraction: if \(M,N\) are terms, then \(MN\) is a
term, and if \(M\) is a term and \(x\) a variable, then \(\lambda x.M\) is a
term.  The metavariables \(M,N,P\) range over this least set.
Application associates to the left and its scope binds more tightly than a
lambda abstraction.  The qualifier \emph{untyped} is necessary because the
grammar contains no type annotations and no typing judgment; it names a
syntactic category, not a class of ill-typed terms.  It differs from the
arrow-only simply typed syntax of Section~\ref{sec:normalization-curry-howard-02}, and that
difference is used by the unrestricted self-application example and the
computability calibration below.
\index{lambda calculus@\(\lambda\)-calculus!untyped}
\index{lambda term!untyped}
\end{hbdefinition}

\begin{hbdefinition}{def:p07-alpha-substitution}{Alpha-equivalence and capture-avoiding substitution}
The free-variable operation on untyped terms is determined by
\[
\begin{aligned}
 \FV(x)&=\{x\},&
 \FV(MN)&=\FV(M)\cup\FV(N),\\
 \FV(\lambda x.M)&=\FV(M)\setminus\{x\}.&&
\end{aligned}
\]
Two terms are \emph{alpha-equivalent} when one is obtained from the other by
finitely many consistent renamings of bound variables to variables fresh for
the corresponding body, within arbitrary term contexts.

For a term \(N\), capture-avoiding substitution \(M[N/x]\) is the recursion
with \(x[N/x]=N\), \(y[N/x]=y\) for \(y\neq x\), and
\((PQ)[N/x]=P[N/x]Q[N/x]\).  At an abstraction it satisfies
\((\lambda x.P)[N/x]=\lambda x.P\); if \(y\neq x\) and
\(y\notin\FV(N)\), then
\((\lambda y.P)[N/x]=\lambda y.(P[N/x])\); otherwise the binder is first
renamed to a variable fresh for \(P,N,x\).  Thus substitution is well defined
up to alpha-equivalence.

The two names are needed to distinguish harmless changes of bound names from
replacement of free occurrences.  Both name metalinguistic relations on
terms, not reduction steps; capture avoidance differs from textual
replacement precisely when a binder would seize a free variable of \(N\).
They are used in the beta contraction and parallel-reduction method below.
\index{alpha-equivalence@\(\alpha\)-equivalence}
\index{substitution!capture-avoiding}
\index{free variable!lambda term@free variable!\(\lambda\)-term}
\end{hbdefinition}

\begin{hbdefinition}{def:p07-beta-conversion}{Beta conversion and normalization}
A \emph{beta redex} is an occurrence of \((\lambda x.M)N\).  One-step beta
reduction is the compatible relation specified on
Sheet~\ref{sheet:lambda-untyped}; that sheet also records capture avoidance and
closure under term contexts.  Write
\(\longrightarrow_{\beta}^{*}\) for its reflexive--transitive closure and
\(M=_{\beta}N\) for the least congruence equivalence relation containing
one-step beta reduction.  A term is in \emph{beta-normal form} when it has no
beta redex.

These names distinguish three standings: beta reduction is a directed
relation on terms, beta conversion is the symmetric congruence it generates,
and normality is a property of one term.  Conversion is not redundant with
reduction because it permits reverse conversion steps, while normality says
nothing yet about whether a starting term can reach a normal form.  The
distinctions are needed by Church--Rosser and its normal-form corollary.

A term is \emph{weakly normalizing} when some finite beta-reduction reaches a
normal form, and \emph{strongly normalizing} when no infinite beta-reduction
starts from it.  These two properties must be named separately because one
asks for a terminating route and the other requires every route to terminate;
they are properties of terms (or, after quantification, of a calculus), not
forms of confluence.  Their nonredundancy is witnessed by the general logical
distinction between existential and universal termination, and the comparison
table uses both.  The self-application term
\((\lambda x.xx)(\lambda x.xx)\) contracts to itself, so the untyped calculus
as a whole has neither weak nor strong normalization.
\index{beta-reduction@\(\beta\)-reduction}
\index{beta-conversion@\(\beta\)-conversion}
\index{normal form!beta@normal form!\(\beta\)}
\index{weak normalization}
\index{strong normalization}
\end{hbdefinition}

\begin{hbmethod}{def:p07-parallel-reduction}{Parallel-reduction method}
Ordinary one-step beta reduction does not itself have the diamond property,
because one step contracts only one chosen redex.  \emph{Parallel reduction}
is the following auxiliary metalinguistic relation on alpha-equivalence
classes of terms.  Write \(M\mathrel{\mathsf{Par}}N\) for the least relation
satisfying four clauses:
\begin{enumerate}
  \item \(x\mathrel{\mathsf{Par}}x\) for every variable \(x\);
  \item if \(M\mathrel{\mathsf{Par}}M'\) and
        \(N\mathrel{\mathsf{Par}}N'\), then
        \(MN\mathrel{\mathsf{Par}}M'N'\);
  \item if \(M\mathrel{\mathsf{Par}}M'\), then
        \(\lambda x.M\mathrel{\mathsf{Par}}\lambda x.M'\);
  \item if \(M\mathrel{\mathsf{Par}}M'\) and
        \(N\mathrel{\mathsf{Par}}N'\), then
        \((\lambda x.M)N\mathrel{\mathsf{Par}}M'[N'/x]\).
\end{enumerate}
Thus one step may contract compatible redexes and their residuals
simultaneously.  Parallel reduction is not a new notion of beta conversion:
its finite transitive closure is proved to be the same as ordinary finite
beta reduction.

The \emph{complete development} operation is the metalinguistic recursion
\[
\begin{aligned}
 \operatorname{dev}(x)&\defeq x,\\
 \operatorname{dev}(\lambda x.M)&\defeq
   \lambda x.\operatorname{dev}(M),\\
 \operatorname{dev}((\lambda x.M)N)&\defeq
   \operatorname{dev}(M)[\operatorname{dev}(N)/x],\\
 \operatorname{dev}(MN)&\defeq
   \operatorname{dev}(M)\operatorname{dev}(N)
   &&\text{when \(M\) is not an abstraction.}
\end{aligned}
\]
The name is needed because this operation contracts all residuals of redexes
already exposed by the starting syntax, rather than choosing one redex or
asserting that arbitrary normalization terminates.  It is an operation on
terms, not a new reduction relation, and its nonredundancy from ordinary
one-step beta reduction is exactly the simultaneous-development property.
That property gives the diamond argument used in
Theorem~\ref{thm:p07-church-rosser}; the annex proves the substitution,
development, simulation, and finite-grid lemmas on which it depends.
\index{parallel reduction}
\index{complete development!lambda calculus@complete development!\(\lambda\)-calculus}
\end{hbmethod}

\begin{hbtheorem}{thm:p07-church-rosser}{proved=annex:p7.lambda.1}{Church--Rosser theorem}
If
\(M\longrightarrow_{\beta}^{*}N_1\) and
\(M\longrightarrow_{\beta}^{*}N_2\), then there is a term \(P\) such that
\(N_1\longrightarrow_{\beta}^{*}P\) and
\(N_2\longrightarrow_{\beta}^{*}P\).
\end{hbtheorem}

The theorem asserts confluence of beta reduction; it does not assert that a
term has a normal form.  Its historical statement and a modern
parallel-reduction development are given in
\cite{ChurchRosser1936,GirardLafontTaylor1989}.

\begin{hbcorollary}{thm:p07-beta-normal-form-unique}{proved=annex:p7.lambda.2}{Uniqueness of beta-normal forms}
If \(M\longrightarrow_{\beta}^{*}N_1\) and
\(M\longrightarrow_{\beta}^{*}N_2\), and both \(N_1,N_2\) are in beta-normal
form, then \(N_1\) and \(N_2\) are alpha-equivalent.
\end{hbcorollary}

\begin{hbwarning}{def:p07-confluence-not-normalization}{Confluence is not normalization}
Theorem~\ref{thm:p07-church-rosser} makes a common descendant available only
after two finite reductions have already been supplied.  It neither produces
a terminating reduction nor excludes an infinite one.  The looping
self-application term above is therefore compatible with Church--Rosser.
\end{hbwarning}

\section[Lambda definability]{Lambda definability and combinatory computation}

\begin{hbdefinition}{def:p07-church-numeral}{Church numeral}
For \(n\in\N\), the \emph{Church numeral for \(n\)} is the closed untyped term
\(\lambda f.\lambda x.f^n x\), where \(f^0x\) abbreviates \(x\) and
\(f^{n+1}x\) abbreviates \(f(f^nx)\).  The name is needed because
lambda-definability below requires one fixed encoding of numerical inputs and
outputs; it names an encoding of a natural number by a term, not a new natural
number.  It differs from the arithmetical numeral notation used in Part~5,
and the fixed encoding is used in Definition~\ref{def:p07-lambda-definable}.
\index{Church numeral}
\end{hbdefinition}

\begin{hbdefinition}{def:p07-lambda-definable}{Lambda-definable partial numerical function}
Let \(f:\N^k\rightharpoonup\N\).  It is \emph{lambda-definable} when there is
a closed untyped term \(F\) such that, for all \(n_1,\ldots,n_k,m\in\N\),
application of \(F\) to the Church numerals for \(n_1,\ldots,n_k\) beta-reduces
to the Church numeral for \(m\) exactly when
\(f(n_1,\ldots,n_k)=m\); when the input is outside \(\dom f\), that application
has no Church numeral as a beta-normal form.  Confluence makes a represented
output unique.

This term is necessary to compare a syntactic computation model with the
closure-theoretic class of Definition~\ref{def:p04-partial-recursive}.  It
names a property of a partial numerical function relative to the fixed Church
encoding, not a property of an arbitrary lambda term.  It differs from
normalization because undefinedness requires only the absence of a numeral
normal form, and it is used in the equivalence theorem immediately below.
\index{lambda-definable function@\(\lambda\)-definable function}
\end{hbdefinition}

\begin{hbtheorem}{thm:p07-lambda-partial-recursive}{cited=Kleene1936}{Lambda-definability and partial recursiveness}
For every arity \(k\), the lambda-definable partial functions
\(\N^k\rightharpoonup\N\) are exactly the partial recursive functions of
Definition~\ref{def:p04-partial-recursive}.
\end{hbtheorem}

Church supplies the lambda-definability framework, Kleene relates the
recursive formalisms, and Church--Rosser supplies the conversion theorem; the
three roles are kept distinct in
\cite{Church1936,Kleene1936,ChurchRosser1936}.

\begin{hbdefinition}{def:p07-sk-combinators}{The S--K combinatory basis and bracket abstraction}
The \emph{S--K combinatory terms} form the least set containing every variable
and the constants \(\mathsf S,\mathsf K\), and closed under application, with
application left-associative.
Its two primitive contractions and their contextual closure are specified only
on Sheet~\ref{sheet:sk-combinatory}.
For one variable \(x\), \emph{bracket abstraction} is the recursive
metalinguistic transformation that sends \(x\) to
\(\mathsf S\mathsf K\mathsf K\), sends a term \(M\) with
\(x\notin\FV(M)\) to \(\mathsf K M\), and sends an application \(MN\) to the
application of \(\mathsf S\) to the abstractions of \(M\) and \(N\).

The fixed basis is needed to state combinatory completeness without leaving
an unspecified stock of primitives.  S--K terms are a syntactic category and
bracket abstraction is a translation on that category; neither is a new
lambda term constructor.  The translation differs from literal term identity
because it removes a binder while preserving application behavior, and that
difference is used in Theorem~\ref{thm:p07-bracket-abstraction}.
\index{combinatory logic!S--K basis}
\index{bracket abstraction}
\end{hbdefinition}

\begin{hbtheorem}{thm:p07-bracket-abstraction}{cited=CurryFeys1958}{Bracket-abstraction property}
For every S--K combinatory term \(M\) and variable \(x\), bracket abstraction
produces a term containing no free occurrence of \(x\) whose application to
\(x\) reduces to \(M\).  Recursively replacing lambda binders by this
transformation gives a combinatory translation of untyped lambda terms; it is
an operational correspondence, not literal identity of the two syntaxes.
\end{hbtheorem}

\begingroup\small
\begin{longtable}{@{}p{.17\textwidth}p{.25\textwidth}p{.25\textwidth}p{.23\textwidth}@{}}
\caption{Untyped lambda and combinatory computation models.}
\label{tab:p07-untyped-computation-models}\\
\toprule
Model & Primitive syntax & Computational step & Expressive calibration \\
\midrule
\endfirsthead
\toprule
Model & Primitive syntax & Computational step & Expressive calibration \\
\midrule
\endhead
Untyped lambda calculus & Variables, application, abstraction &
Capture-avoiding beta contraction & Church-numeral functions are exactly the
partial recursive numerical functions
(Theorem~\ref{thm:p07-lambda-partial-recursive}). \\
S--K combinatory logic & Variables, application, \(\mathsf S,\mathsf K\) &
The two contractions on Sheet~\ref{sheet:sk-combinatory} & Bracket
abstraction eliminates binders; completeness is cited from
\cite{CurryFeys1958}. \\
\bottomrule
\end{longtable}
\endgroup

\section{System F: explicit polymorphism}

\begin{hbdefinition}{def:p07-system-f}{Church-style System F}
The \emph{Church-style System F} used here fixes disjoint countably infinite
sets of type variables \(a,b,\ldots\) and term variables
\(x,y,\ldots\).  Its types form the least set containing every type variable
and closed under the function-type former
\(\alpha\mathbin{\to_{\mathrm{ty}}}\beta\) and universal type formation
\(\forall a.\alpha\); \(\alpha,\beta\) range over that set.  Its terms form
the least set containing every term variable and closed under typed
abstraction \(\lambda x{:}\alpha.M\), application \(MN\), type abstraction
\(\Lambda a.M\), and type application \(M[\alpha]\); \(M,N\) range over that
set.  Its judgments have the form
\(\Delta_{\mathrm{ty}};\Theta_{\mathrm{ty}}\derives_{\mathsf F}M{:}\alpha\).
The type-variable context \(\Delta_{\mathrm{ty}}\) lists distinct type
variables; the term-variable context \(\Theta_{\mathrm{ty}}\) assigns
distinct term variables types well formed over that type context.  The exact
formation, typing, freshness, and capture-avoidance conditions are on
Sheet~\ref{sheet:system-f}.

The style qualifier is necessary because erasing type annotations, type
abstractions, and type applications gives a different assignment problem.
This definition names an explicitly annotated proof-term calculus, not all
formulations called System~F.  It differs from STLC by quantification over
types, and the distinction controls both the encodings and the decision rows
below.
\index{System F!Church style}
\index{polymorphic lambda calculus@polymorphic \(\lambda\)-calculus}
\end{hbdefinition}

\begin{hbdefinition}{def:p07-system-f-impredicativity}{Impredicativity in System F}
System~F type quantification is \emph{impredicative}: in an instance of
\(\forall a.\alpha\), the variable \(a\) may be replaced by any well-formed
System~F type, including one that itself contains universal type
quantification.  This names a property of the type-formation and elimination
regime, not a term-reduction rule.  It differs from mere polymorphic notation
because the quantified collection is not stratified below the quantified
type, and it is used by the encodings in the next example.
\index{impredicativity!System F}
\end{hbdefinition}

\begin{hbexample}{def:p07-system-f-encodings}{Boolean and numeral encodings}
The types
\[
  \forall a.(a\mathbin{\to_{\mathrm{ty}}}a
    \mathbin{\to_{\mathrm{ty}}}a)
  \qquad\text{and}\qquad
  \forall a.((a\mathbin{\to_{\mathrm{ty}}}a)
    \mathbin{\to_{\mathrm{ty}}}a\mathbin{\to_{\mathrm{ty}}}a)
\]
support the familiar Church encodings of Booleans and natural numbers.  They
are examples of impredicative encoding, not a claim that every desired
extensional law for every datatype is internally derivable.
\end{hbexample}

\begin{hbtheorem}{thm:p07-system-f-confluence}{cited=GirardLafontTaylor1989}{Confluence of System F reduction}
For the Church-style terms and the combined compatible reduction generated by
term-beta and type-beta contraction on Sheet~\ref{sheet:system-f}, any two
finite reductions from the same term have a common reduct.
\end{hbtheorem}

\begin{hbtheorem}{thm:p07-system-f-strong-normalization}{cited=Girard1972}{Strong normalization for System F}
Every term typable on Sheet~\ref{sheet:system-f} admits no infinite sequence
of term-beta and type-beta contractions.
\end{hbtheorem}

Under Curry--Howard, the types of System~F are formulas of second-order
intuitionistic propositional logic, universal type abstraction is
second-order universal introduction, and typed terms record derivations.  The
claim is calculus-relative: it uses the rules on Sheet~\ref{sheet:system-f},
just as Definition~\ref{def:p07-curry-howard} fixes the STLC correspondence.

\begin{hbdefinition}{def:p07-system-f-presentation-styles}{Church and Curry presentations}
For the explicitly annotated syntax of Definition~\ref{def:p07-system-f},
\emph{checking} asks whether a displayed judgment
\(\Delta_{\mathrm{ty}};\Theta_{\mathrm{ty}}\derives_{\mathsf F}M{:}\alpha\) follows from
the sheet.  In the \emph{Curry-style assignment} presentation, terms contain
neither term-binder type annotations nor explicit type abstraction and
application; type checking asks, for a given erased \(M\) and given
\(\alpha\), whether \(M\) can be assigned \(\alpha\), while typability asks
whether some type can be assigned.

These style names are needed because the input data of the two decision
problems differ.  They name two formal presentations and their metalinguistic
decision problems, not two spellings of one algorithm.  Erasure connects but
does not identify them, and the distinction is used by the next two theorems
and Table~\ref{tab:p07-lambda-calculi-comparison}.
\index{System F!Church style!type checking}
\index{System F!Curry style}
\index{typability!System F}
\end{hbdefinition}

\begin{hbtheorem}{thm:p07-system-f-church-checking}{cited=GirardLafontTaylor1989}{Church-style checking}
Given finite well-formed \(\Delta_{\mathrm{ty}},\Theta_{\mathrm{ty}}\), an
explicit Church-style term \(M\), and a type \(\alpha\), it is decidable
whether the judgment
\[
 \Delta_{\mathrm{ty}};\Theta_{\mathrm{ty}}
 \derives_{\mathsf F}M{:}\alpha
\]
is derivable on Sheet~\ref{sheet:system-f}.
\end{hbtheorem}

\begin{hbtheorem}{thm:p07-system-f-curry-undecidable}{cited=Wells1999}{Undecidability for Curry-style System F}
In the Curry-style System~F assignment presentation fixed in
Definition~\ref{def:p07-system-f-presentation-styles}, there are effective
reductions in both directions between typability and type checking, and both
decision problems are undecidable.
\end{hbtheorem}

\begin{proofboundary}{def:p07-system-f-parametricity}{Parametricity pointer}
Reynolds's relational parametricity interprets a polymorphic type by a
relation-preservation condition and proves an abstraction theorem for the
selected relational model \cite{Reynolds1983}.  This chapter records that
connection only: it does not define the logical relation, prove the
abstraction theorem, or infer unrestricted ``free theorems'' from the two
example encodings.  Reynolds's earlier account of polymorphism is
\cite{Reynolds1974}.
\index{parametricity}
\end{proofboundary}

\section[Intensional type theory]{An intensional Martin--L\"of boundary and the CIC pointer}

\begin{proofboundary}{def:p07-mltt-boundary}{Selected MLTT calculus boundary}
The literature name ``Martin--L\"of type theory'' covers inequivalent choices
of judgments, equality, universes, and type formers. This handbook selects the
exact core on \hyperref[sheet:mltt-intensional]{\texttt{sheet:mltt-intensional}}:
context, type, term, type-equality, and term-equality judgments; explicit
weakening, substitution, equivalence, congruence, and conversion rules;
dependent products, dependent sums, intensional identity types; and a
cumulative predicative Russell universe hierarchy. This is a formal calculus,
not a name for the whole literature family. Its nonredundant exclusions are
equality reflection, general recursion, quotients, proof irrelevance,
univalence, and higher inductive types. The selected boundary is used by the
propositions-as-types entry and by the narrower normalization comparison below
\cite{MartinLof1984,NordstromPeterssonSmith1990,Hofmann1997}.
\index{Martin-Lof type theory@Martin--L\"of type theory}
\index{dependent product}
\index{dependent sum}
\index{identity type!intensional}
\index{universe!type theory}
\end{proofboundary}

The sheet fixes syntax and derivability, not a normalization theorem for every
listed former.  In particular, the theorem below concerns a proper subcalculus
with dependent functions and one Russell universe; it is not evidence for the
sheet's dependent sums, identity types, or cumulative hierarchy.  Normalization
and decidable type checking remain calculus-specific metatheorems
\cite{MartinLof1984,AbelCoquandDybjer2007}.

\begin{hbdefinition}{def:p07-propositions-as-types}{Propositions as types and proof relevance}
For this entry, \emph{propositions as types} means that a proposition is
represented by a type and a proof by an inhabitant of that type.  This is an
interpretive correspondence between propositions and types, not the axiom
that every two inhabitants are equal.  \emph{Proof relevance} is the
possibility that inhabitants, including inhabitants of identity types, remain
distinguishable unless a separate proof-irrelevance principle is assumed.
The second name is necessary to block the common but invalid inference from
propositions-as-types to proof irrelevance; it names a possible feature of
the judgmental theory and is used to delimit the MLTT row of the comparison
table.
\index{propositions as types}
\index{proof relevance}
\index{proof irrelevance}
\end{hbdefinition}

\begin{hbtheorem}{thm:p07-mltt-normalization-decision}{cited=AbelCoquandDybjer2007}{Normalization and typed equality in the Abel--Coquand--Dybjer subcalculus}
For the intensional Martin--L\"of core with typed equality judgments,
dependent functions, and one Russell-style universe specified by Abel,
Coquand, and Dybjer, normalization by evaluation computes normal forms for
well-typed terms and decides typed beta-eta equality.  Consequently fully
annotated type checking is decidable for that same syntax.  This statement
does not cover dependent sums, identity types, cumulative or stratified
universe hierarchies, or the full selected calculus in
\ref{def:p07-mltt-boundary}.
\end{hbtheorem}

\begin{hbtheorem}{thm:p07-coc-strong-normalization}{cited=CoquandHuet1988}{Strong normalization for the Calculus of Constructions}
Every well-typed term of the exact Calculus of Constructions presented by
Coquand and Huet is strongly normalizing.  This statement is about that
calculus, not about a later inductive extension.
\end{hbtheorem}

\begin{proofboundary}{def:p07-cic-pointer}{CIC pointer}
The \emph{Calculus of Inductive Constructions} (CIC) is recorded here only as
the literature name for an inductive extension of the dependent Calculus of
Constructions.  Coquand--Huet \cite{CoquandHuet1988} supplies the predecessor
calculus and the theorem immediately above; it does not by itself define or
establish the metatheory of CIC.  The phrase here is therefore a
metalinguistic pointer, not the definition of a handbook calculus: no exact
inductive declaration, positivity checker, conversion relation, cumulative
universe hierarchy, or constraint set has been fixed.  CIC consequently has
no normalization or decision claim in this chapter.
\index{Calculus of Constructions}
\index{Calculus of Inductive Constructions}
\end{proofboundary}

\begin{proofboundary}{def:p07-hott-cubical-boundary}{HoTT, univalence, and cubical type theory boundary}
This chapter gives only a pointer to homotopy type theory and univalent
foundations; it does not add
univalence, a homotopical semantics, or the identity-type hierarchy to the
MLTT boundary above.  It likewise leaves higher
inductive types and cubical composition/computation outside scope.  These are
scope boundaries, not compressed technical entries, and no theorem in
this chapter depends on either extension.
\index{homotopy type theory}
\index{univalence}
\index{cubical type theory}
\index{higher inductive type}
\end{proofboundary}

\begin{proofboundary}{def:p07-proof-nets-goi-boundary}{Proof-net and GoI boundary}
The proof-net discussion is a pointer only \cite{Girard1987}: no correctness
criterion, sequentialization theorem, or net-based cut elimination is
developed.  Geometry of interaction is likewise only a pointer
with no operator-algebraic execution model or GoI interpretation imported.
These boundaries are adjacent to the linear
logic material elsewhere in Part~7, not additional lambda-calculus results.
\index{proof net}
\index{geometry of interaction}
\end{proofboundary}

\section{Comparison table}

Table~\ref{tab:p07-lambda-calculi-comparison} separates reduction facts from
typing decision problems.  ``Confluent'' always refers to the reduction
relation fixed by the row; ``strongly normalizing'' quantifies only over the
terms specified by the row.  Sources for the table are
\cite{Barendregt1984,GirardLafontTaylor1989,Girard1972,Wells1999,
AbelCoquandDybjer2007}.

\begin{landscape}
\begingroup\scriptsize
\setlength{\tabcolsep}{3pt}
\begin{longtable}{@{}p{.95in}p{1.30in}p{1.45in}p{1.55in}p{1.55in}@{}}
\caption{Lambda calculi and type theories compared on confluence,
normalization, checking, and typability.}
\label{tab:p07-lambda-calculi-comparison}\\
\toprule
System / presentation & Reduction and confluence & Normalization & Checking &
Typability / inference \\
\midrule
\endfirsthead
\toprule
System / presentation & Reduction and confluence & Normalization & Checking &
Typability / inference \\
\midrule
\endhead
Untyped lambda calculus & Beta reduction is confluent
(Theorem~\ref{thm:p07-church-rosser}). & Not every term is weakly or strongly
normalizing; the self-application loop is a counterexample. & No intrinsic
typing judgment, hence no native type-checking problem. & No native
typability problem. \\
STLC, explicit binders & Beta reduction on well-typed terms is confluent. &
Every well-typed term is strongly normalizing by the source theorem in
Section~\ref{sec:normalization-curry-howard-02}. & Syntax-directed checking
of a fully annotated term is decidable. & For the arrow-only simple-type
assignment problem, constraint generation and first-order unification decide
typability; this is presentation-specific, not a System~F claim. \\
System F, Church style & Combined term-beta and type-beta reduction is
confluent (Theorem~\ref{thm:p07-system-f-confluence}). & Every well-typed term
is strongly normalizing
(Theorem~\ref{thm:p07-system-f-strong-normalization}). & Explicit checking is
decidable (Theorem~\ref{thm:p07-system-f-church-checking}). & Erased-term
typability is not the input problem of this presentation. \\
System F, Curry style & Underlying beta reduction is confluent; typing is a
separate assignment relation. & Every typable term is strongly normalizing by
the System~F normalization theorem. & Given an erased term and a proposed
type, checking is undecidable (Wells). & Typability is equivalent to checking
for decidability and is undecidable
(Theorem~\ref{thm:p07-system-f-curry-undecidable}). \\
Abel--Coquand--Dybjer MLTT core: dependent functions and one Russell universe &
The row tracks its typed beta-eta judgmental equality; no distinct blanket
rewrite-confluence claim is added. & Normalization by evaluation computes
normal forms for its well-typed terms; extensions require separate theorems. &
Typed beta-eta equality and fully annotated checking are decidable for that
exact calculus
(Theorem~\ref{thm:p07-mltt-normalization-decision}). & Dependent-type
inference is not claimed decidable by this entry. \\
CoC / CIC boundary & CoC has its cited reduction theory; no CIC reduction
relation is fixed here. & CoC strong normalization is
Theorem~\ref{thm:p07-coc-strong-normalization}; CIC has no blanket cell. & No
CIC claim. & No CIC claim. \\
\bottomrule
\end{longtable}
\endgroup
\end{landscape}
